% ===================== Preface =====================
\chapter*{Preface}
\addcontentsline{toc}{chapter}{Preface}
\markboth{Preface}{Preface}

\section*{Why Salam exists}

Almost every programming language in wide use today was designed around
English. Its keywords are English words \code{if}, \code{while},
\code{return} and its tutorials, error messages, and documentation assume
that the reader is comfortable reading English. For hundreds of millions of
people whose first language is Persian or Arabic, that assumption is a wall to
climb before the real learning can even begin.

\textbf{Salam} (from the word \fa{سلام}, meaning \emph{peace}) was created to
lower that wall. It is, at its heart, an ordinary modern programming language:
statically typed, compiled, fast, with structs, generics, interfaces, closures,
a real standard library, and the ability to call C. What makes it unusual is
that it is \emph{bilingual}. The very same language can be written with English
keywords or with Persian keywords, and both produce identical programs. A loop
is spelled \code{while} in one and \fa{تاوقتی} in the other; a function is
\code{func} or \fa{تابع}. The compiler, \code{salam}, understands both.

This means a student in Tehran can write their first program in their own
language, reading their own script, with error messages in their own
words and then, when they are ready, switch the very same concepts into
English to work with the wider software world. The ideas transfer perfectly,
because they are \emph{the same ideas}.

\section*{Who this book is for}

This book assumes \emph{no prior programming experience}. If you have never
written a line of code in your life, start at Chapter~1 and walk forward one
step at a time every term is explained the first time it appears, and every
concept is shown with a small, complete program you can run yourself.

If you are already a programmer coming from C, Python, Go, or Rust, you will
find Salam comfortably familiar. You can move quickly through the early
chapters, treat the \emph{In Depth} boxes as your main course, and use the
appendices as a quick reference. Salam will feel like a language you already
half-know.

The book is therefore written on two levels at once:

\begin{itemize}[leftmargin=1.4em]
  \item The \textbf{main text} is friendly and gradual. It uses plain language,
        small examples, and analogies, and it never assumes you have seen a
        concept before.
  \item The \textbf{In Depth} boxes dig into how things actually work the
        memory model, monomorphization, the C that Salam generates,
        performance trade-offs. They are clearly marked as optional. Skip them
        on a first reading; come back when you are curious.
\end{itemize}

\section*{How to use this book}

Read with your hands, not just your eyes. Programming is a craft, and crafts are
learned by doing. Every chapter contains complete programs type them in, run
them, and then \emph{change} them. Break them on purpose and see what error the
compiler gives. The single most effective way to learn a language is to argue
with its compiler until you both agree.

Each chapter ends with exercises. Attempt them before peeking at
Appendix~D, where solutions to selected exercises are collected. The boxes
sprinkled through the text mean:

\begin{itemize}[leftmargin=1.4em]
  \item \textbf{Note} a clarification or a fact worth remembering.
  \item \textbf{Tip} a practical shortcut or good habit.
  \item \textbf{Watch out} a common mistake; read these twice.
  \item \textbf{In Depth} optional, deeper material for the curious.
\end{itemize}

\section*{A note on scope}

Salam also ships with a built-in \emph{layout} dialect for describing web pages
(it compiles small \code{layout:} blocks straight to HTML and CSS). That side of
the language is a topic of its own and is \emph{not} covered here. This book is
about Salam as a general-purpose programming language: variables, types, control
flow, functions, data structures, the standard library, concurrency, and the
foreign function interface. Everything you learn here is pure Salam
programming.

\begin{note}
All the example programs in this book are drawn from, or directly modelled on,
the real, tested examples that ship with the Salam compiler (under
\code{examples/en/} and \code{tests/en/}). They are not pseudo-code; they
compile and run.
\end{note}

\section*{Conventions}

Salam source code is shown in framed, syntax-highlighted boxes with line
numbers:

\begin{salam}
func main:
    println "Hello, Salam!"
end
\end{salam}

\noindent Terminal commands and program output are shown in plain grey boxes:

\begin{progout}
$ salam exec hello.salam
Hello, Salam!
\end{progout}

\noindent Inline code such as \code{func} or a filename like
\code{hello.salam} appears in a monospaced font. When a Persian keyword is shown
beside its English form, it appears like this: \code{func}~/~\code{\fa{تابع}}.

\bigskip
\noindent Welcome to Salam. Let us begin.

\vspace{1em}
\hfill The Salam Language Team
\clearpage
